%% toc: P15 172 172.11,172.12
\section{Proofs for robust conclusions, causal effects and robust decisions}
\label{sec:annex-p15-w02b-p15-070}

Notation is that of
\cref{def:p15-w02b-p15-070-robust-conclusion,def:p15-w02b-p15-070-asymmetric-family,def:p15-w02b-p15-070-three-claims,def:p15-w02b-p15-070-confounded-shift}.
Throughout, \(Z\) denotes a standard normal variable when it appears
inside \(m_k\), and \(m_k(c)=E_\Phi\psi_k(Z-c)\) as in
\cref{thm:p15-huber-asymmetric-bias}.

\begin{annexitem}{annex:p15-w02b-p15-070-reduction}{thm:p15-w02b-p15-070-reduction}{Reduction to signs of differences}
\begin{proof}
\emph{Item 1.} If \(\inf_\Theta D>0\), then
\(D(\theta)\geq\inf_\Theta D>0\) for every member. If every
\(D(\theta)>0\), each is a lower bound candidate exceeding zero, so the
infimum, the greatest lower bound, is at least zero. For
\(\Theta=(0,1]\) and \(D(\theta)=\theta\), each value is positive and
\(\inf_\Theta D=0\).

Let \(\Theta\) be a nonempty compact topological space and \(D\) lower
semicontinuous, and put \(m=\inf_\Theta D\in[-\infty,\infty)\).
For real \(r>m\), the set \(K_r=\{\theta:D(\theta)\leq r\}\) is closed
by lower semicontinuity and nonempty by the definition of the infimum.
If \(r_1<\cdots<r_p\) are real numbers exceeding \(m\), then
\(K_{r_1}\cap\cdots\cap K_{r_p}=K_{r_1}\neq\varnothing\).
Compactness and this finite intersection property give a point
\(\theta_0\) in \(\bigcap_{r>m}K_r\). Then \(D(\theta_0)\leq r\) for
every real \(r>m\), hence \(D(\theta_0)\leq m\); since \(D(\theta_0)\)
is real and at least \(m\), the infimum is finite and equals
\(D(\theta_0)\). Thus it is a minimum. If \(\min D>0\), the conclusion
is robust by the first sentence; if it is robust, the value
\(D(\theta_0)=\min D\) is positive.

\emph{Item 2.} Established means \(D(\theta)>0\) for all members and
reversed means \(D(\theta)<0\) for all members. Because \(\Theta\) is
nonempty, both cannot hold. The verdict is unresolved exactly when
neither holds, that is, when some member has \(D(\theta)\leq0\) and some
member has \(D(\theta')\geq0\). A member with \(D=0\) serves for both, and
two members of opposite strict sign serve for one each.

\emph{Item 3.} For real numbers, the chain
\(y_1>y_2>\cdots>y_m\) holds exactly when \(y_j>y_{j+1}\) for every
\(j<m\): the chain contains each adjacent inequality, and conversely,
for \(j<l\), transitivity of \(>\) applied to \(y_j>y_{j+1}>\cdots>y_l\)
gives \(y_j>y_l\). Apply this at each member to
\(y_j=q_{\sigma(j)}(\theta)\). The ranking holds at every member exactly
when each adjacent difference is positive at every member.

\emph{Item 4.} If \(\inf_\Theta q_i>\sup_\Theta q_j\), then
\(q_i(\theta)\geq\inf q_i>\sup q_j\geq q_j(\theta)\). In the example,
\(q_1-q_2=1\) at every member, while the ranges are \([1,11]\) and
\([0,10]\), which overlap.

For the product family, \(\inf q_i\in[-\infty,\infty)\) and
\(\sup q_j\in(-\infty,\infty]\) because the functions are real valued on
nonempty sets, so \(\inf q_i-\sup q_j\in[-\infty,\infty)\) is defined.
At every member \((\theta',\theta'')\),
\(q_i(\theta')-q_j(\theta'')\geq\inf q_i-\sup q_j\). Conversely choose
\(\theta'_p\) with \(q_i(\theta'_p)\) decreasing to \(\inf q_i\) and
\(\theta''_p\) with \(q_j(\theta''_p)\) increasing to \(\sup q_j\) (both
exist by the definitions of infimum and supremum, including the infinite
cases). Because the coordinates can be chosen independently,
\((\theta'_p,\theta''_p)\in\Theta\), and
\(q_i(\theta'_p)-q_j(\theta''_p)\) tends to \(\inf q_i-\sup q_j\).
Hence the infimum of the difference equals that value.
\end{proof}
\end{annexitem}

\begin{annexitem}{annex:p15-w02b-p15-070-margins}{thm:p15-w02b-p15-070-margins}{Additive margins}
\begin{proof}
\emph{Item 1.} At every member,
\(D(\theta)\geq d_0-\sum_r|e_r(\theta)|\geq d_0-E\). If \(d_0>E\), this
is positive. Symmetrically
\(D(\theta)\leq d_0+E<0\) when \(d_0<-E\).

\emph{Item 2.} At a member \(\theta^\dagger\) with
\(e_r(\theta^\dagger)=-E_r\) for every \(r\),
\(D(\theta^\dagger)=d_0-E\leq0\), so \(D>0\) fails there and is not
robust. In the product case choose, for each \(r\), a coordinate
\(\theta_r\in\Theta_r\) with \(e_r(\theta_r)=-E_r\). The tuple
\((\theta_1,\ldots,\theta_R)\) belongs to \(\Theta\), and each
\(e_r\), depending only on its own coordinate, equals \(-E_r\) there.

\emph{Item 3.} At \(\theta_1\), \(e_1+e_2=E_1-E_1=0\); at
\(\theta_2\), \(e_1+e_2=-E_1+E_1=0\). With \(E_2=E_1\), the bounds hold,
\(D=d_0\) at both members, and \(D>0\) is robust for every \(d_0>0\),
although item 1 certifies it only for \(d_0>2E_1\).
\end{proof}
\end{annexitem}

\begin{annexitem}{annex:p15-w02b-p15-070-linear-ranking}{thm:p15-w02b-p15-070-linear-ranking}{Linear variance shares}
\begin{proof}
Fix a member, write \(\mu_j=EX_j\), and fix \(i\). Independence of
\(X_i\) from \(X_{-i}=(X_j)_{j\ne i}\) gives
\(E[X_j\mid X_i]=\mu_j\) for \(j\ne i\), so
\[
 E[Y\mid X_i]=a_0+a_iX_i+\sum_{j\ne i}a_j\mu_j
\]
almost surely, and \(\operatorname{Var}(E[Y\mid X_i])=a_i^2v_i\).
Likewise \(E[Y\mid X_{-i}]=a_0+a_i\mu_i+\sum_{j\ne i}a_jX_j\), so
\(Y-E[Y\mid X_{-i}]=a_i(X_i-\mu_i)\). This variable is independent of
\(X_{-i}\), hence its conditional second moment given \(X_{-i}\) equals
its unconditional second moment \(a_i^2v_i\) almost surely, and
\(E[\operatorname{Var}(Y\mid X_{-i})]=a_i^2v_i\).
Independence also gives \(V=\operatorname{Var}(Y)=\sum_ja_j^2v_j\).
Dividing by \(V>0\) gives \(S_i=S_{T_i}=a_i^2v_i/V\), as in
\cref{def:p15-global-indices}.

For \(i\ne j\), \(S_i-S_j=(a_i^2v_i-a_j^2v_j)/V\), whose sign is that
of the numerator. Over the box, \(v_i\) and \(v_j\) vary independently,
so by item 4 of \cref{thm:p15-w02b-p15-070-reduction} the numerator has
infimum \(a_i^2\ell_i-a_j^2u_j\), attained at any point of the box with
\(v_i=\ell_i\) and \(v_j=u_j\). The strict ranking holds at every member
exactly when this attained minimum is positive, and the weak ranking
exactly when it is nonnegative.

For the share threshold put \(x=a_i^2v_i\) and
\(y=\sum_{j\ne i}a_j^2v_j\), so \(S_i=x/(x+y)\) with \(x,y\geq0\),
\(x+y=V>0\). For fixed \(y\) the map \(x\mapsto x/(x+y)\) has derivative
\(y/(x+y)^2\geq0\) where \(x+y>0\); for fixed \(x\) the map
\(y\mapsto x/(x+y)\) has derivative \(-x/(x+y)^2\leq0\). Let \(\bar v\)
be the corner with \(\bar v_i=\ell_i\) and \(\bar v_j=u_j\) for
\(j\ne i\); it lies in the box, so its value of \(V\), which is the
displayed denominator, is positive. For any point \(v\) of the box,
first lower \(v_i\) to \(\ell_i\): the intermediate point
\((\ell_i,v_{-i})\) lies in the box, so its \(V\) is positive and the
share does not increase. Then raise each \(v_j\), \(j\ne i\), to
\(u_j\), again staying in the box; the share does not increase. Hence
the minimum share over the box is attained at \(\bar v\) and equals the
displayed fraction. The threshold conclusion holds throughout exactly
when this minimum exceeds \(s_0\).

With \(a=(1,2)\), \(a_1^2\ell_1\geq a_2^2u_2\) reads
\(\ell_1\geq4u_2\).
\end{proof}
\end{annexitem}

\begin{annexitem}{annex:p15-w02b-p15-070-contamination-margin}{thm:p15-w02b-p15-070-contamination-margin}{The contamination margin}
\begin{proof}
\emph{Properties of \(m_k\).} The score \(\psi_k\) is continuous,
nondecreasing, odd, bounded by \(k\) in absolute value, and
1-Lipschitz.
(a) Continuity: if \(c_p\to c\), then \(\psi_k(Z-c_p)\to\psi_k(Z-c)\)
pointwise with bound \(k\), so \(m_k(c_p)\to m_k(c)\) by dominated
convergence.
(b) Oddness: \(Z\) and \(-Z\) have the same law, so
\(m_k(-c)=E\psi_k(-Z+c)=-E\psi_k(Z-c)=-m_k(c)\); in particular
\(m_k(0)=0\).
(c) Strict decrease: for \(c<c'\),
\(\psi_k(z-c)-\psi_k(z-c')\geq0\) for every \(z\), and it is positive
when \(z-c'<k\) and \(z-c>-k\), that is, on the interval
\((c-k,c'+k)\), which has positive standard normal probability.
Hence \(m_k(c)>m_k(c')\).
(d) Range: \(\psi_k(Z-c)>-k\) on the event \(\{Z>c-k\}\) of positive
probability and \(\psi_k\geq-k\) everywhere, so \(m_k(c)>-k\); by (b),
\(m_k(c)<k\). As \(c\to\infty\), \(\psi_k(Z-c)\to-k\) pointwise with
bound \(k\), so \(m_k(c)\to-k\).
Consequently \(m_k\) maps \([0,\infty)\) continuously and strictly
decreasingly onto \((-k,0]\). For \(0\leq\varepsilon<1/2\) the number
\(-\varepsilon k/(1-\varepsilon)\) lies in \((-k,0]\), so \(c_*\geq0\),
the solution of \(m_k(c_*)=-\varepsilon k/(1-\varepsilon)\) in
\cref{thm:p15-huber-asymmetric-bias}, exists and is unique, and by (b),
\(m_k(-c_*)=\varepsilon k/(1-\varepsilon)\).

\emph{The population score of a member.} For
\(F=(1-\varepsilon)\Phi+\varepsilon H\) let
\(\lambda_F(s)=E_F\psi_k(X-s)\). Since \(\psi_k\) is bounded,
\[
 \lambda_F(s)=(1-\varepsilon)m_k(s)+\varepsilon\int\psi_k(x-s)\,H(dx).
\]
The second term is nonincreasing in \(s\) and the first is strictly
decreasing when \(\varepsilon<1\), so \(\lambda_F\) is strictly
decreasing and has at most one zero. By
\cref{thm:p15-huber-asymmetric-bias} it has exactly one, \(c(F)\).

\emph{Item 1.} The clipped estimate is the midpoint of a minimizing
interval, which is translation equivariant by
\cref{thm:p15-convex-location-minimizers}. With \(X_i=\xi+W_i\),
\(T_{n,k}(X_1,\ldots,X_n)=\xi+T_{n,k}(W_1,\ldots,W_n)\), and the second
term converges almost surely to \(c(F)\) by
\cref{thm:p15-huber-asymmetric-bias}.

\emph{Item 2, inclusion.} Let \(0\leq\varepsilon<1/2\). For
\(s>c_*\), the bound \(\psi_k\leq k\) and (c) give
\[
 \lambda_F(s)\leq(1-\varepsilon)m_k(s)+\varepsilon k
 <(1-\varepsilon)m_k(c_*)+\varepsilon k=0,
\]
and for \(s<-c_*\), using \(\psi_k\geq-k\),
\(\lambda_F(s)\geq(1-\varepsilon)m_k(s)-\varepsilon k
>(1-\varepsilon)m_k(-c_*)-\varepsilon k=0\). Hence
\(|c(F)|\leq c_*\) for every member. If
\(\xi\in\Xi_\varepsilon(t)\), then \(t-\xi=c(F)\in[-c_*,c_*]\).

\emph{Item 2, every point is attained.} If \(\varepsilon=0\), then
\(c_*=0\), the only member is \(\Phi\), and \(c(\Phi)=0\) because
\(m_k(0)=0\); thus \(\Xi_0(t)=\{t\}\). Let \(0<\varepsilon<1/2\) and
\(d\in[-c_*,c_*]\). Fix \(M>c_*+k\) and put
\[
 u=\frac12\Bigl(1-\frac{(1-\varepsilon)m_k(d)}{\varepsilon k}\Bigr),
 \qquad H=u\,\delta_M+(1-u)\,\delta_{-M}.
\]
Since \(m_k(d)\in[m_k(c_*),m_k(-c_*)]\) by (c), the ratio
\((1-\varepsilon)m_k(d)/(\varepsilon k)\) lies in \([-1,1]\), so
\(u\in[0,1]\) and \(H\) is a probability. Because \(|d|\leq c_*\),
\(M-d>k\) and \(-M-d<-k\), so \(\psi_k(M-d)=k\) and
\(\psi_k(-M-d)=-k\). Therefore
\[
 \lambda_F(d)=(1-\varepsilon)m_k(d)+\varepsilon k(2u-1)
 =(1-\varepsilon)m_k(d)-(1-\varepsilon)m_k(d)=0,
\]
and \(c(F)=d\) by uniqueness of the zero. For any real \(t\), the
member \(F\) with \(\xi=t-d\) gives \(t(\xi,F)=t\). Thus every point of
\([t-c_*,t+c_*]\) belongs to \(\Xi_\varepsilon(t)\), each realized by a
contaminant on the two points \(\pm M\).

\emph{Item 3.} By item 2, the consistent centers form
\([t-c_*,t+c_*]\), a compact interval. The function
\(\xi\mapsto\xi-\xi_0\) is continuous, so by item 1 of
\cref{thm:p15-w02b-p15-070-reduction} \(\xi>\xi_0\) holds on it exactly
when its minimum \(t-c_*-\xi_0\) is positive; likewise \(\xi<\xi_0\)
holds exactly when the maximum \(t+c_*-\xi_0\) is negative; otherwise
the verdict is unresolved by item 2 of that theorem.

For two separately contaminated samples, the consistent pairs
\((\xi_A,\xi_B)\) form the product
\([t_A-c_*,t_A+c_*]\times[t_B-c_*,t_B+c_*]\), since each sample's
contaminant is chosen separately. By item 4 of
\cref{thm:p15-w02b-p15-070-reduction}, the infimum of \(\xi_A-\xi_B\) is
\(t_A-c_*-(t_B+c_*)\), attained at a corner. The ranking is robust exactly
when this is positive. If both samples share one law \(F\), item 1
gives \(t_A-t_B=(\xi_A+c(F))-(\xi_B+c(F))=\xi_A-\xi_B\).

\emph{Item 4.} The map \(g(\varepsilon)=\varepsilon k/(1-\varepsilon)\)
is continuous and strictly increasing on \([0,1/2)\), with
\(g(0)=0\) and \(g(\varepsilon)\to k\) as \(\varepsilon\uparrow1/2\).
Let \(\mu\) be the inverse of the restriction of \(m_k\) to
\([0,\infty)\); it is continuous and strictly decreasing from
\((-k,0]\) onto \([0,\infty)\), because a continuous strictly monotone
bijection between intervals has a continuous inverse. Then
\(c_*(\varepsilon)=\mu(-g(\varepsilon))\) is continuous and strictly
increasing, and \(c_*(0)=\mu(0)=0\). For every \(C>0\),
\(m_k(C)>-k\) by (d); once \(\varepsilon\) is so close to \(1/2\) that
\(-g(\varepsilon)<m_k(C)\), strict decrease of \(m_k\) gives
\(c_*(\varepsilon)>C\). Hence \(c_*(\varepsilon)\to\infty\) as
\(\varepsilon\uparrow1/2\).

Let \(\delta>0\). By item 3, \(\xi>\xi_0\) is robust over
\(\mathcal F_\varepsilon\) exactly when \(c_*(\varepsilon)<\delta\).
Because \(m_k\) is strictly decreasing on \([0,\infty)\), this is
equivalent to \(m_k(c_*(\varepsilon))>m_k(\delta)\), that is,
\(-g(\varepsilon)>m_k(\delta)\), or
\[
 \frac{\varepsilon}{1-\varepsilon}<a,\qquad a=\frac{-m_k(\delta)}{k}.
\]
By (c) and (d), \(-k<m_k(\delta)<0\), so \(0<a<1\). The map
\(\varepsilon\mapsto\varepsilon/(1-\varepsilon)\) is a strictly
increasing bijection of \([0,1)\) onto \([0,\infty)\) with inverse
\(y\mapsto y/(1+y)\), so the condition reads
\(\varepsilon<a/(1+a)=-m_k(\delta)/(k-m_k(\delta))=\varepsilon^*(\delta)\),
and \(0<\varepsilon^*(\delta)<1/2\) because \(0<a<1\).
The families are nested: for \(\varepsilon\leq\varepsilon'<1\),
\((1-\varepsilon)\Phi+\varepsilon H=(1-\varepsilon')\Phi+\varepsilon'H'\)
with the probability
\(H'=((\varepsilon'-\varepsilon)\Phi+\varepsilon H)/\varepsilon'\)
when \(\varepsilon'>0\). Hence \(\varepsilon^*(\delta)\) is the
robustness radius of the conclusion in this nested collection.

\emph{Item 5.} Let \(1/2\leq\varepsilon<1\) and \(M>k\). For
\(H=\delta_M\) and any \(s<M-k\),
\[
 \lambda_F(s)=(1-\varepsilon)m_k(s)+\varepsilon k
 >-(1-\varepsilon)k+\varepsilon k\geq0,
\]
by (d). Since \(\lambda_F\) is strictly decreasing with zero \(c(F)\),
\(c(F)\geq M-k\). Reflection gives the opposite direction: if \(F'\) is
the law of \(-X\) for \(X\) with law \(F\), then
\(F'=(1-\varepsilon)\Phi+\varepsilon H'\) with \(H'\) the reflection of
\(H\), and \(\lambda_{F'}(s)=E\psi_k(-X-s)=-\lambda_F(-s)\), whose unique
zero is \(-c(F)\). Thus for each real \(r\) there are members with
\(c\geq r\) and with \(c\leq r\). Given \(t\) and \(\xi_0\), a member
with \(c(F)\geq t-\xi_0\) gives the consistent center
\(\xi=t-c(F)\leq\xi_0\), and a member with \(c(F)\leq t-\xi_0\) gives a
consistent center \(\xi\geq\xi_0\). Neither strict conclusion is robust.
\end{proof}
\end{annexitem}

\begin{annexitem}{annex:p15-w02b-p15-070-linearized-margin}{thm:p15-w02b-p15-070-linearized-margin}{A strict lower bound}
\begin{proof}
\emph{Derivative of \(m_k\).} For \(h\ne0\),
\[
 \frac{\psi_k(z-c-h)-\psi_k(z-c)}{h}\in[-1,0],
\]
because \(\psi_k\) is nondecreasing and 1-Lipschitz. For \(z\) with
\(|z-c|\ne k\), \(\psi_k\) is differentiable at \(z-c\) with derivative
\(\mathbf 1\{|z-c|<k\}\), so the quotient tends to
\(-\mathbf 1\{|z-c|<k\}\) as \(h\to0\). The excluded set
\(\{c-k,c+k\}\) has normal probability zero. By dominated convergence,
\[
 m_k'(c)=-P(|Z-c|<k)=-\bigl(\Phi(c+k)-\Phi(c-k)\bigr)=-g(c),
\]
which is continuous in \(c\). Hence \(m_k(c)=-\int_0^cg(s)\,ds\) by the
fundamental theorem of calculus and \(m_k(0)=0\).

\emph{Comparison with the slope at zero.} The function \(g\) has
derivative \(\phi(s+k)-\phi(s-k)\). For \(s>0\) and \(k>0\),
\((s+k)^2-(s-k)^2=4sk>0\), so \(\phi(s+k)<\phi(s-k)\) and \(g\) is
strictly decreasing on \([0,\infty)\). Thus \(g(s)<g(0)=A(k)\) for
\(s>0\), and for \(c>0\)
\[
 m_k(c)=-\int_0^cg(s)\,ds>-A(k)c,
\]
the inequality being strict because the continuous integrand is below
\(A(k)\) on \((0,c]\).

\emph{The bound.} For \(0<\varepsilon<1/2\),
\(m_k(c_*)=-\varepsilon k/(1-\varepsilon)<0=m_k(0)\), so \(c_*>0\).
Applying the last inequality at \(c=c_*\),
\[
 -\frac{\varepsilon k}{1-\varepsilon}=m_k(c_*)>-A(k)c_*,
 \qquad\text{so}\qquad
 c_*>\frac{\varepsilon k}{(1-\varepsilon)A(k)}.
\]

\emph{The consequence.} If the margin satisfies
\(\varepsilon k/((1-\varepsilon)A(k))<\delta<c_*\), then
\(\xi_0=t-\delta\) lies in \((t-c_*,t)\subseteq\Xi_\varepsilon(t)\) by
item 2 of \cref{thm:p15-w02b-p15-070-contamination-margin}, so a member
consistent with \(t\) has center exactly \(\xi_0\), and \(\xi>\xi_0\) is
not robust; yet the linearized value is smaller than \(\delta\).

\emph{The numbers.} Evaluate the closed form of \(m_k\) given in
\cref{thm:p15-huber-asymmetric-bias}. For \(\varepsilon=0.1\) and
\(k=1.14\), the target is
\(-\varepsilon k/(1-\varepsilon)=-0.12666\ldots\), and to seven decimals
\[
\begin{aligned}
 m_{1.14}(0.1702)&=-0.1265312,\\
 m_{1.14}(0.1703)&=-0.1266050,\\
 m_{1.14}(0.1704)&=-0.1266789.
\end{aligned}
\]
The first two exceed the target and the third is below it, so by strict
decrease \(0.1703<c_*<0.1704\), and \(c_*>0.1702\). Also
\(A(1.14)=0.7457137\) to seven decimals, so the linearized value is
\(0.114/(0.9\times0.7457137)=0.16986\) to five decimals. The margin
\(0.1702\) exceeds both \(0.16986\) and \(0.17\) and is less than
\(c_*\). For the two other table entries, the targets
\(-0.0196970\) (\(\varepsilon=0.01\), \(k=1.95\)) and
\(-0.00264264\) (\(\varepsilon=0.001\), \(k=2.64\)), to the digits
shown, satisfy
\[
\begin{aligned}
 m_{1.95}(0.020755)&=-0.0196925>-0.0196970,\\
 m_{1.95}(0.020765)&=-0.0197020<-0.0196970,\\
 m_{2.64}(0.0026645)&=-0.00264241>-0.00264264,\\
 m_{2.64}(0.0026650)&=-0.00264291<-0.00264264,
\end{aligned}
\]
which give the stated intervals for \(c_*\).
\end{proof}
\end{annexitem}

\begin{annexitem}{annex:p15-w02b-p15-070-decisions}{thm:p15-w02b-p15-070-decisions}{Robust decisions, regret and minimax}
\begin{proof}
\emph{Item 1.} Regret is nonnegative, so \(\bar r(a)=0\) exactly when
\(r(a,\theta)=0\) for every \(\theta\), that is, when \(a\) is optimal at
every member. If such an action exists, the smallest maximal regret is
zero, and the actions attaining it are those with \(\bar r=0\), the
robustly optimal actions.

\emph{Item 2.} If \(a\) is robustly optimal, then
\(L(a,\theta)\leq L(b,\theta)\) for every \(b\) and \(\theta\), so
\(\sup_\theta L(a,\theta)\leq\sup_\theta L(b,\theta)\) for every \(b\);
\(a\) is minimax. In the first example \(L_*=(0,5)\), so \(a\) is
optimal at both members, and both actions have supremum \(5\), so \(b\)
is minimax too, although it is not optimal at \(\theta_1\). In the
second example \(L_*=(0,0)\): the third action is optimal at neither
member, the first is not optimal at \(\theta_2\), and the second is not
optimal at \(\theta_1\), so no action is robustly optimal. The suprema
are \(10,10,6\), so the third action is the unique minimax action.

\emph{Item 3.} Let \(L'(a,\theta)=\lambda L(a,\theta)+h(\theta)\) with
\(\lambda>0\). At each \(\theta\), the map \(x\mapsto\lambda x+h(\theta)\)
is strictly increasing, so the minimizers of \(L'(\cdot,\theta)\) are
those of \(L(\cdot,\theta)\), and
\(L'_*(\theta)=\lambda L_*(\theta)+h(\theta)\). Hence the robustly optimal
actions agree, \(r'=\lambda r\), \(\bar r'=\lambda\bar r\), and the
minimax-regret actions agree. In the example the suprema of
\((-2,1)\) and \((0,0)\) are \(1\) and \(0\), so abstaining is the unique
minimax action. After adding \(h=(2,-1)\) the loss vectors are
\((0,0)\) and \((2,-1)\), with suprema \(0\) and \(2\), and acting is the
unique minimax action.

\emph{Item 4.} By hypothesis, acting is optimal at \(\theta\) exactly
when \(q(\theta)\geq c\), so it is optimal at every member exactly when
\(q\geq c\) on \(\Theta\); likewise for abstaining with \(q\leq c\).
A robustly optimal action therefore exists exactly when \(q\geq c\)
throughout or \(q\leq c\) throughout, which fails exactly when some member
has \(q>c\) and another has \(q<c\). Acting is the unique optimal action
at \(\theta\) exactly when acting is optimal and abstaining is not, that
is, when \(q(\theta)\geq c\) and not \(q(\theta)\leq c\), which is
\(q(\theta)>c\). This holds at every member exactly when \(q>c\) is
robust. For the displayed loss, acting is optimal at \(\theta\) exactly
when \(c-q(\theta)\leq0\) and abstaining exactly when
\(0\leq c-q(\theta)\), which is the hypothesis.
\end{proof}
\end{annexitem}

\begin{annexitem}{annex:p15-w02b-p15-070-minimax-not-robust}{thm:p15-w02b-p15-070-minimax-not-robust}{Minimax without optimality at each member}
\begin{proof}
\emph{Membership.} With \(H=\Phi\), which is symmetric about zero,
\((1-\varepsilon)\Phi+\varepsilon\Phi=\Phi\in\mathcal C_\varepsilon\).

\emph{The mean at \(\Phi\).} At center zero,
\(\sqrt n\,\overline X_n\) has the standard normal law for every \(n\).
Hence \(E_\Phi\min(n\overline X_n^2,C^2)=E\min(Z^2,C^2)\) does not
depend on \(n\), and as \(C\to\infty\) it increases to \(EZ^2=1\) by
monotone convergence. Thus \(\mathcal R(\overline X;\Phi)=1\).

\emph{The clipped estimate at \(\Phi\).} By
\cref{thm:p15-huber-normal-variance} with \(F=\Phi\),
\(\sqrt nT_{n,k}\Rightarrow N(0,V_k(\Phi))\). For fixed \(C\), the map
\(x\mapsto\min(x^2,C^2)\) is bounded and continuous, so
\(E_\Phi\min(nT_{n,k}^2,C^2)\to E\min(V_k(\Phi)Z^2,C^2)\); the
\(\liminf\) is this limit. As \(C\to\infty\) it increases to
\(V_k(\Phi)\). Thus \(\mathcal R(T_{\cdot,k};\Phi)=V_k(\Phi)\).
Because \(\Phi\) has no atoms, \(A_\Phi(k)=\Phi((-k,k))=A(k)\). Since
\(\phi'(x)=-x\phi(x)\), integration by parts gives
\[
 \int_{-k}^kx^2\phi(x)\,dx=\bigl[-x\phi(x)\bigr]_{-k}^{k}
 +\int_{-k}^k\phi(x)\,dx=A(k)-2k\phi(k),
\]
and \(P(|X|>k)=2(1-\Phi(k))\), so
\(E_\Phi\min(X^2,k^2)=A(k)-2k\phi(k)+2k^2(1-\Phi(k))=B(k)\). Hence
\(V_k(\Phi)=B(k)/A(k)^2\).

\emph{The strict inequality.} Using \(\int_k^\infty x\phi(x)\,dx=\phi(k)\),
\[
\begin{aligned}
 E_\Phi[X\psi_k(X)]&=\int_{-k}^kx^2\phi(x)\,dx+2k\int_k^\infty x\phi(x)\,dx\\
 &=A(k)-2k\phi(k)+2k\phi(k)=A(k),
\end{aligned}
\]
and \(E_\Phi\psi_k(X)^2=E_\Phi\min(X^2,k^2)=B(k)\). By the
Cauchy--Schwarz inequality,
\(A(k)^2\leq E_\Phi X^2\cdot E_\Phi\psi_k(X)^2=B(k)\). Equality would
require \(\psi_k(X)=\lambda X\) almost surely for a constant \(\lambda\),
because \(E_\Phi X^2=1>0\). On the event \(\{0<X<k\}\), which has positive
probability, this forces \(\lambda=1\); on \(\{X>k\}\), also of positive
probability, \(\psi_k(X)=k<X\). So the inequality is strict and
\(B(k)/A(k)^2>1\). For \(k=1\), \(B(1)=0.5160586\) and
\(A(1)^2=0.4660649\) to seven decimals, with ratio \(1.1073\) to four
decimals.

\emph{Worst cases.} For \(M>0\), the symmetric contaminant
\(H=(\delta_M+\delta_{-M})/2\) gives a member \(F\) with mean zero and
variance \(1-\varepsilon+\varepsilon M^2\), as in
\cref{thm:p15-huber-location-failure-cases}. By
\cref{thm:p15-clt-iid}, \(\sqrt n\,\overline X_n\) converges in law to
a centered normal law with this variance, and the argument used for
\(T_{n,k}\) above gives \(\mathcal R(\overline X;F)=1-\varepsilon
+\varepsilon M^2\). As \(M\to\infty\) this is unbounded, so
\(\mathcal R_\varepsilon(\overline X)=\infty\) for \(\varepsilon>0\).
For the normalized \(k\),
\(\mathcal R_\varepsilon(T_{\cdot,k})=V^*(k,\varepsilon)=1/I_*\) by
\cref{thm:p15-huber-location-minimax}, which is finite.
\end{proof}
\end{annexitem}

\begin{annexitem}{annex:p15-w02b-p15-070-consistency-identification}{thm:p15-w02b-p15-070-consistency-identification}{Consistency forces identification}
\begin{proof}
\emph{Item 1.} Under \(M\), the observations
\((O_1,\ldots,O_n)\) have the product law \(P_M^{\otimes n}\) on
\(\mathcal O^n\), and the law of \(T_n\) is the image of that product
law under \(T_n\). If \(P_{M_1}=P_{M_2}\), the product laws coincide for
every \(n\), and so do the image laws. A property defined from the
sequence of these laws, or from the laws of the procedure under the
members of a neighborhood that is itself determined by \(P_M\), is a
function of \(P_M\) alone and therefore takes the same value at the two
models.

\emph{Item 2.} Suppose \(P_{M_1}=P_{M_2}=P\) and put
\(q_i=Q(M_i)\). By item 1, the law \(\Lambda_n\) of \(T_n\) is the same
under both models, and the hypothesis gives
\(\Lambda_n(\{y:|y-q_1|>\eta\})\to0\) and
\(\Lambda_n(\{y:|y-q_2|>\eta\})\to0\) for every \(\eta>0\).
If \(q_1\ne q_2\), take \(\eta=|q_1-q_2|/3\). The sets
\(\{|y-q_1|\leq\eta\}\) and \(\{|y-q_2|\leq\eta\}\) are disjoint, and
for large \(n\) each has \(\Lambda_n\)-probability above \(1/2\), which
is impossible for a probability measure. Hence \(q_1=q_2\), and \(Q\) is
identified.

\emph{Item 3.} This is the contrapositive of item 2. Each class of
\cref{thm:p16-causal-nonidentification} contains two models with equal
observational laws and unequal target values, so the target is not
identified there.
\end{proof}
\end{annexitem}

\begin{annexitem}{annex:p15-w02b-p15-070-separation}{thm:p15-w02b-p15-070-separation}{The three claims in the confounded shift model}
\begin{proof}
\emph{Item 1.} The assignment \(X=Z\) makes \(X\) a fair binary
variable. On \(\{X=x\}\) we have \(Z=x\), so
\(Y=\beta x+\gamma x+W=\alpha x+W\); since \(W\) is independent of
\(Z=X\), the conditional law of \(Y\) given \(X=x\) is the law of
\(\alpha x+W\). The joint law of \((X,Y)\) is therefore determined by
\(\alpha\) and \(F\).

Under \(\operatorname{do}(X=x)\) the assignments of \(Z\), \(W\) and
\(Y\) are kept, so \(Y^{X=x}=\beta x+\gamma Z+W\) at every outcome, and
\(Y^{X=1}=Y^{X=0}+\beta\) pointwise. Thus \(\beta\) is a shift effect.
It is the only one: if \(Y^{X=0}+s\) and \(Y^{X=0}+s'\) had the same law
with \(d=s'-s\ne0\), the characteristic function \(\chi\) of
\(Y^{X=0}\) would satisfy \(\chi(u)=e^{iud}\chi(u)\) for all real
\(u\). Since \(\chi\) is continuous with \(\chi(0)=1\), it is nonzero
for \(|u|\) small, and for small \(u\ne0\) the factor \(e^{iud}\) differs
from one, a contradiction.

The models \(M(\beta,\gamma,F)\) and \(M(\beta+\gamma,0,F)\) have the
same \(\alpha\) and \(F\), hence the same observational law, and shift
effects \(\beta\) and \(\beta+\gamma\). When \(\gamma\ne0\) these
differ, which is nonidentification.

\emph{Item 2.} With uncontaminated records and \(F\in\mathcal
F_\varepsilon\), the observations in arm \(x\) are independent with the
law of \(\alpha x+W\). By item 1 of
\cref{thm:p15-w02b-p15-070-contamination-margin}, the arm estimate
converges almost surely to \(\alpha x+c(F)\), and the contrast to
\(\alpha\), whatever the contaminant. This equals \(\beta\) exactly when
\(\gamma=0\). With separate record contaminations, the arm laws are
\(F_x(\cdot-\alpha x)\) and the limits are \(\alpha x+c(F_x)\), so the
contrast converges to \(\tau=\alpha+c(F_1)-c(F_0)\). By item 2 of
\cref{thm:p15-w02b-p15-070-contamination-margin}, each \(c(F_x)\) lies
in \([-c_*,c_*]\) and every value there is realized by some
\(F_x\in\mathcal F_\varepsilon\); the two are chosen separately, so
\(\tau-\alpha\) fills \([-2c_*,2c_*]\).

\emph{Item 3.} In the subclass \(\gamma=0\), \(F=\Phi\), the law of
\(Y\) given \(X=1\) is \(\Phi(\cdot-\beta)\), whose mean is \(\beta\);
equal observational laws therefore give equal \(\beta\), which is
identification, and \(\beta=\alpha\) is the shift effect by item 1. The
arm means are \(0\) and \(\beta\), so their difference is \(\beta\). The
population score of \(\Phi\) vanishes at zero because \(m_k(0)=0\), so
\(c(\Phi)=0\) and the clipped contrast tends to \(\beta\).

Contaminate the records of arm \(X=1\) alone, so that its law is
\(F_1(\cdot-\beta)\) with \(F_1=(1-\varepsilon)\Phi+\varepsilon H\),
\(0<\varepsilon<1/2\), while arm \(X=0\) keeps \(\Phi\). The clipped
contrast tends to \(\beta+c(F_1)-c(\Phi)=\beta+c(F_1)\), within \(c_*\)
of \(\beta\) by item 2 of
\cref{thm:p15-w02b-p15-070-contamination-margin}. If \(H\) has a mean
\(\mu_H\), the contaminated arm has mean
\(\beta+(1-\varepsilon)\cdot0+\varepsilon\mu_H\), so the difference of
arm means is \(\beta+\varepsilon\mu_H\); taking \(H=\delta_{\mu}\) with
\(\mu\) arbitrary shows that it is unbounded over the family. If
\(\int|x|\,H(dx)=\infty\), then
\(\int|x|\,F_1(dx)\geq\varepsilon\int|x|\,H(dx)=\infty\), so the
contaminated arm has no mean and the difference of means is undefined.

\emph{Item 4.} With uncontaminated records, item 2 gives the contrast limit
\(\beta+\gamma\) at each member of \(\Theta_\Gamma\). The members
consistent with an observed limit \(\alpha\) are those with
\(\beta=\alpha-\gamma\), \(\gamma\in[0,\Gamma]\); their shift effects
fill \([\alpha-\Gamma,\alpha]\), with each value attained. The
conclusion that the shift effect exceeds \(c\) holds at all of them
exactly when the attained minimum \(\alpha-\Gamma\) exceeds \(c\). By
item 4 of \cref{thm:p15-w02b-p15-070-decisions}, applied to the
consistent members with \(q\) the shift effect, adoption is then the
unique optimal action at every consistent member. The consistent effects
form a single point exactly when \(\Gamma=0\); for \(\Gamma>0\) two
consistent members have the same observational law and different
effects, so the effect is not identified among them.

With separate record contaminations, a member is determined by
\(\gamma\in[0,\Gamma]\), \(F_0,F_1\in\mathcal F_\varepsilon\) and
\(\beta\), with contrast limit \(\beta+\gamma+c(F_1)-c(F_0)\). Given an
observed limit \(\tau\), the consistent effects are
\(\beta=\tau-\gamma-c(F_1)+c(F_0)\), where \(\gamma\), \(c(F_1)\) and
\(c(F_0)\) range separately over \([0,\Gamma]\), \([-c_*,c_*]\) and
\([-c_*,c_*]\), all values attained by item 2 of
\cref{thm:p15-w02b-p15-070-contamination-margin}. By item 4 of
\cref{thm:p15-w02b-p15-070-reduction}, the effects fill
\([\tau-\Gamma-2c_*,\tau+2c_*]\), and the conclusion that the effect
exceeds \(c\) is robust exactly when \(\tau-2c_*-\Gamma>c\).

\emph{Item 5.} The model \(M(\alpha-\gamma,\gamma,F)\) has
\(\beta+\gamma=\alpha\), hence by item 1 the same observational law as
any other model with these \(\alpha\) and \(F\), and its shift effect is
\(\alpha-\gamma<c\) when \(\gamma>\alpha-c\). If also \(\gamma>\Gamma\),
it is not a member of \(\Theta_\Gamma\), so the verdict of item 4 says
nothing about it, and adoption is not optimal for it.
\end{proof}
\end{annexitem}
